\documentclass[a4paper,11pt]{letter}
\pagestyle{empty}

\usepackage{amsmath}
\usepackage{amsfonts}
\usepackage{amssymb}
\usepackage{graphicx}
\usepackage[margin=2.54cm]{geometry}

\begin{document}

% Whoami

% Brief description of why you want to pursue research studies

% About what your academic interests are.

% How they relate to your previous studies and future goals.


Dear Hiring Committee,

My name is Simon Sj\"oling and I am applying for the doctoral student position in Autonomous Systems at KTH Royal Institute of Technology. I am currently finishing my master's thesis on real-time neural soft shadow rendering, which will be completed in June 2026, concluding five years of study at Chalmers University of Technology (M.Sc. in Data Science and AI and B.Sc. in Electrical Engineering, Civilingenjör).

% Brief description of why you want to pursue research studies

% Answer:
% - I have always liked the research process: reading literature, setting up experiments, conducting experiments and writing in detail about it.
% - I want to be a part of the frontier of the research, both contributing but also to learn

Even before starting university education at Chalmers, I had my goal set on a research degree.
% The parts of my studies I have enjoyed the most have been the ones closest to research: reading literature, formulating a problem, implementing a solution, evaluating the results and writing clearly about the work.
% The research process of reading literature, designing experiments or projects, evaluating the results and scientific writing has always been the most enjoyable part of my studies.
The parts of my studies that I have found most enjoyable are the ones closest to research, such as reading scientific literature, formulating a problem, implementing a solution, evaluating the results and writing clearly about the work.
I have found that I work well in projects where there is room to go deeper and understand why something works rather than only making it work. Pursuing doctoral studies feels like a natural continuation of that way of working.

% About what your academic interests are and how they relate to your previous studies and future goals.

% Answer: 
% - Machine learning 
% - Physical applications of ML (robotics), not just in software
% - Mathematics

During my master's studies, my primary academic interest has been machine learning, both the mathematical and algorithmic theory behind the methods and their use in real engineering problems. I have especially enjoyed courses in probabilistic machine learning, AI systems, stochastic processes, computer vision, high performance computing and nonlinear optimisation. In my thesis, I explore how small specialized neural networks can approximate computationally expensive physical phenomena such as soft shadows for computer graphics, with preliminary results achieving up to 3 times faster computation times.
% Although the application domain is graphics, the work has strengthened skills that I believe transfer well to autonomous systems research: efficient implementation, careful experimental evaluation and reasoning about accuracy and computational cost under real-time constraints.
Although the application domain of my thesis is in computer graphics, I believe that many of the skills I have obtained transfer to autonomous systems, as both are real-time systems that have unique constraint of accuracy and computational cost.

My electrical engineering background has also given me exposure to automatic control, linear systems, signal processing, mechatronics and low-level implementation. Together with my interest in machine learning and mathematics, this makes autonomous systems a direction I am very motivated to move toward. Formal methods are a newer direction for me,
% but one I am interested in developing,
% especially because they connect mathematical specification with reliable behavior in real systems.
but one I am very willing to develop, especially since the use mathematics to model reliability of autonomous systems seems really interesting.

Outside of school, I have spent a lot of time building software projects, including Deep Q-learning and MuZero agents, simple C++ and CUDA-based machine learning frameworks and various game engine projects. These projects have strengthened both my programming ability and challenged me to understand real research papers. Through my work at Ericsson and Chalmers Teknologkonsulter, I have also gained experience developing software in professional settings, working in teams and taking responsibility for real technical deliverables.

I am very excited at the prospect of contributing to the research of the group at KTH and learning from experts in the field.
% My goal is to develop into a researcher who can combine mathematical methods, machine learning and strong implementation skills to build reliable autonomous systems.
I believe my background gives me a strong foundation for this project and I am highly motivated to develop further through doctoral research with the group.

% I am strongly interested in machine learning, especially when it comes to the engineering aspect of applying the technology to solve real problems. In my thesis, I explore how neural networks can approximate computationally expensive physical phenomena such as lighting and shadows for computer graphics, with an emphasis on accuracy and efficiency under strict real-time constraints ($> 60 \, \text{Hz}$).
% A central part of this work has been iterating on the model and its implementation, using evaluations and ablation studies to identify which design choices have the highest impact on accuracy and performance. My main goal after my thesis is to contribute to the field of machine learning through research. Reading about this position made me excited to apply because improving the quality of brain MSI can have real impact on the understanding of brain processes and the opportunity to develop new generative methods for super-resolution and imputation closely aligns with my goals.

% My studies have given me a solid foundation in mathematics, physics and computer science. During my master's, I have particularly focused on AI through coursework and project work which has included implementing, optimizing and evaluating machine learning systems. I have worked with deep learning for vision, transformers and generative methods such as diffusion models and I am comfortable with modern ML workflows and libraries. Many of the courses I have taken, specifically image analysis and computer vision, have included building models specifically for medical applications, such as classifying medical imaging. I have found these projects to be particularly interesting, given the importance of the applications. Additionally, I have taken a course in biomedical engineering and have a basic knowledge of medical imaging.

% Much of my coursework has involved group projects, which I have enjoyed and during my work at Ericsson I gained additional experience working in larger teams, taking responsibility for deliverables and presenting results.

% Beyond machine learning, I have also chosen courses mainly in mathematics, such as Non-linear Optimization, Partial Differential Equations, Stochastic Processes and Bayesian Inference and financial mathematics, which have strengthened my mathematical understanding and advanced problem-solving skills. I have also taken High Performance Computing, which has been valuable for running computationally demanding tasks efficiently on both CPUs and, importantly for ML, GPUs.

% Outside of school, I have always had a passion for tinkering with software, which has resulted in numerous software implementations of RL agents, simple C++ and CUDA-based ML frameworks and more. This has strengthened my practical programming skills and my interest in the lower-level implementation of machine learning algorithms.

% I am very excited at the prospect of contributing to the research of your group at KTH and learning from experts in the field. I believe that my background could be a good starting point and I am highly motivated to develop through doctoral research with this group.

Best regards,\\
Simon Sjöling

\end{document}
